%% Szkielet wygenerowany z docs/02_spis_tresci.md (plan v1, 11.11.2026).
%% Blok planroboczy usuń, gdy podrozdział jest napisany.

\chapter{Dyskusja}
\label{ch:dyskusja}

% objętość planowana: ok. 3,5 s.

\section{Progi agenta LLM a~model uczenia motywowanego}
\label{sec:6-1}

\begin{planroboczy}[ok. 1,5 s.]
\begin{itemize}
  \item Czy agent LLM zachowuje się jak agent z~progiem?
  \item Co odpowiada WTA, a~co bólom abstrakcyjnym.
  \item Modele progów odpowiedzi z~etologii (podział pracy u~owadów społecznych): ta sama sigmoidalna funkcja progowa z~parametrem stromości; progi stałe i~progi wzmacniane doświadczeniem jako analogia do kalibracji progu promptem (PB2); krytyka: same progi nie wyjaśniają całego zachowania (\cite{bonabeau1996quantitative}, \cite{theraulaz1998response}, \cite{ulrich2021response}).
\end{itemize}

\zrodlaplanowane{\cite{starzyk2017needs, zawislak2025budowa, graham2015opportunistic, keramati2014homeostatic, bonabeau1996quantitative, theraulaz1998response, ulrich2021response, yoshida2025linking}}
\end{planroboczy}

\section{Wyniki na tle literatury o~wrażliwości na prompt}
\label{sec:6-2}

\begin{planroboczy}[ok. 1 s.]
\begin{itemize}
  \item Porównanie skali efektów z~wcześniejszymi badaniami wrażliwości na prompt i~z~progami przełączenia z~badania punkty–ból.
  \item Wyniki z~preprintów oznaczamy w~tekście jako wyniki preprintu (\texttt{03}, zasada 4): \cite{keeling2024can}, \cite{masumori2025survival}, \cite{tagliabue2026pain}, \cite{schlatter2026incomplete}, \cite{mikaelson2025beyond}.
\end{itemize}

\zrodlaplanowane{\cite{sclar2024quantifying, keeling2024can, masumori2025survival, tagliabue2026pain, schlatter2026incomplete, mizrahi2024state, sharma2024towards, cedro2026cash, mikaelson2025beyond}}
\end{planroboczy}

\section{Ograniczenia i~kierunki dalszych badań}
\label{sec:6-3}

\begin{planroboczy}[ok. 1 s.]
\begin{itemize}
  \item Jeden model, uproszczone środowisko, dyskretne kategorie stanu, budżet prób.
  \item Ryzyko antropomorfizacji.
  \item Propozycje dla kolejnych prac, w~tym styk z~pracą Kingi.
  \item Zależność od preprintów: wątek z~3.4 opiera się głównie na preprintach (stan X.2026); w~tekście zaznaczamy, które twierdzenia ich dotyczą.
  \item Wyniki z~preprintów oznaczamy w~tekście jako wyniki preprintu (\texttt{03}, zasada 4): \cite{ren2026ai}.
\end{itemize}

\zrodlaplanowane{\cite{sharkey2025could, li2024evaluating, ren2026ai, shanahan2023role, hagendorff2023machine, atil2025nondeterminism, dominguezolmedo2024questioning, song2025good, butlin2026identifying, long2024taking}}
\end{planroboczy}
